\bookchapter{第六十一章：Swift 快速上手——蘋果開發嘅第一語言}

\chapterintro{呢一章我哋正式學語言。唔使緊張——Swift 係出名「對新手友好」，而且你唔需要背語法書：先掌握 20\% 嘅語法，就已經睇得明同寫得出 80\% 嘅日常代碼。淨低嘅一邊做項目一邊查。記住我哋嘅學習原則：睇得明 > 記得住，識用 > 精通。}

\booksection{一、變數同類型：畀數據改個名}

程式就係「數據 + 邏輯」。變數，就係畀數據改個名儲起佢。Swift 用 let 同 var 嚟聲明：

\begin{lstlisting}[language=Swift]
let name = "張三"        // let：常量（推薦預設用呢個）
var age = 18             // var：變數（之後仲可以改）

age = 19                 // <i class="fas fa-check-circle mr-1"></i> var 可以改
// name = "李四"         // <i class="fas fa-times-circle mr-1"></i> let 唔可以改，編譯報錯

let score: Int = 100     // 都可以明確寫明類型
let pi = 3.14            // 唔寫都會自動推斷

print("你好，\(name)")   // 字串插值：反斜槓 + 括號
\end{lstlisting}

常用類型一張表記住：

\begin{center}
\begin{tabularx}{\linewidth}{|X|X|X|}
\hline
\textbf{類型} & \textbf{存乜嘢} & \textbf{示例} \tabularnewline
\hline
Int & 整數 & 42 \tabularnewline
Double & 小數 & 3.14 \tabularnewline
String & 文字 & "你好" \tabularnewline
Bool & 真/假 & true / false \tabularnewline
{}[Int]（陣列） & 一組數據 & {}[1, 2, 3] \tabularnewline
\hline
\end{tabularx}
\end{center}

\begin{center}\small\color{gray}Swift 常用基礎類型\end{center}

\begin{tcbinfo}{小貼士：let 優先原則}
Swift 社區嘅鐵律：預設用 let，真係有需要先用 var。可以用 let 嘅地方用 var，Xcode 會出黃字警告你。呢個習慣幫你避開一大類「變數畀人意外改咗」嘅 bug——代碼變化越少，就越容易推理。

\end{tcbinfo}

\booksection{二、可選值（Optional）：Swift 嘅“防崩護身符”}

呢個係 Swift 最重要嘅概念，冇之一。一句講晒：當一個值可能“唔存在”嗰陣，Swift 會迫你處理佢——喺語言層面消滅咗“空指針崩潰”呢個行業噩夢。

\begin{lstlisting}[language=Swift]
var name: String = "張三"   // 普通類型：絕對唔會係 nil
// name = nil                // <i class="fas fa-times-circle mr-1"></i> 編譯錯誤，直接擋住

var nickname: String? = nil  // 加個問號：允許爲 nil（可選值）
nickname = "三兒"

// 用可選值，必須"解包"：
if let n = nickname {        // if let：安全解包
    print("你好，\(n)")
} else {
    print("仲未有暱稱")
}

let len = nickname?.count ?? 0   // ?? 畀預設值（空合運算符）
print(len)
\end{lstlisting}

\begin{tcbwarn}{注意：! 係"高危操作"}
代碼入面寫 nickname!（強制解包）等於對編譯器嗌“我保證佢唔係 nil”——一估錯，App 就會即場崩潰。新手期原則：盡量用 if let / ??，唔好用 !。Xcode 補全代碼嗰陣有時會順手加個 !，記得改返佢。

\end{tcbwarn}

\booksection{三、函數同閉包：將邏輯打包成“按鈕”}

函數係一段有名嘅邏輯，好似搖控器上面嘅按鈕。閉包（Closure）就係“冇名嘅函數”，可以隨手遞畀人：

\begin{lstlisting}[language=Swift]
// 函數：名字喺前，參數同返回類型喺後
func greet(name: String) -> String {
    return "你好，\(name)！"
}
print(greet(name: "張三"))     // 調用時要帶參數名

// 閉包：就地定義，輕快直接
let numbers = [1, 2, 3, 4]
let doubled = numbers.map { $0 * 2 }   // $0 = 當前元素
// [2, 4, 6, 8]

let big = numbers.filter { $0 > 2 }    // 篩選
// [3, 4]
\end{lstlisting}

SwiftUI 入面閉包無處不在：按鈕嘅 Button(“點我”) \{ ... \} 後面嗰對大括號就係一個閉包。map、filter、sorted 呢幾招“集合連招”請記實——將來 90\% 嘅列表處理都係佢哋嘅組合。

\booksection{四、結構體同協議：Swift 嘅“樂高”體系}

Swift 官方嘅偏好好明確：優先用結構體（struct），要到需要繼承嗰陣先用類（class）。結構體係“值類型”——複製嗰陣各自各嘅，天然少一類 bug：

\begin{lstlisting}[language=Swift]
// 結構體：裝數據嘅盒
struct Record {
    var id = UUID()
    var amount: Double
    var category: String
    var date: Date

    // 結構體自帶"逐一成員構造器"：直接就建到
}
let r = Record(amount: 25.5, category: "餐飲", date: Date())

// 協議：約定"必須具備邊啲能力"
protocol Identifiable {
    var id: UUID { get }
}
struct Record: Identifiable { ... }   // 遵守協議 = 承諾實現其要求
\end{lstlisting}

SwiftUI 入面幾乎所有嘢都係 struct（View 就係協議）。呢句而家唔明唔緊要，第 62 章會日日見面——喺 Swift 嘅世界入面，先擁抱 struct。

\booksection{五、寫畀零基礎嘅語法速查卡}

\begin{center}
\begin{tabularx}{\linewidth}{|X|X|}
\hline
\textbf{想做乜嘢} & \textbf{寫法} \tabularnewline
\hline
定義常量／變數 & \texttt{let x = 1}/\texttt{var x = 1} \tabularnewline
字串插值 & "你好，\textbackslash{}(name)" \tabularnewline
定義函數 & func 名字(參數: 類型) -> 返回類型 \{ ... \} \tabularnewline
判斷 & if 條件 \{ ... \} else \{ ... \} \tabularnewline
遍歷 & for x in 陣列 \{ ... \} \tabularnewline
安全解包 & if let x = 可選值 \{ ... \} \tabularnewline
預設值 & x ?? 預設值 \tabularnewline
定義結構體 & \texttt{struct X \{ var a: Int; var b: String \}} \tabularnewline
閉包 & \texttt{\{ \$0 * 2 \}} \tabularnewline
\hline
\end{tabularx}
\end{center}

\begin{center}\small\color{gray}Swift 語法速查卡（新手版）\end{center}

\begin{tcbpractice}{嚟動手：喺 Playgrounds 玩 15 分鐘 Swift}
有 Mac 就開 Xcode Playground；冇 Mac 就用 iPad 上嘅 Swift Playgrounds 或者網頁版演練場。將本章每一段代碼都行一次，改改數字同文字，故意整幾個錯誤睇吓 Xcode 點提示。呢個“先玩後學”嘅 15 分鐘，會令你第 62 章上手嗰陣輕鬆一倍。

\end{tcbpractice}

\begin{tcbsummary}{本章小結}
\begin{itemize}
  \item let 優先，var 慎用；類型自動推斷，冒號後面顯式聲明都得
  \item 可選值係 Swift 嘅靈魂：可能係空就必須處理，if let / ?? 係護身符，! 係高危操作
  \item 函數調用要帶參數名；閉包 = 冇名嘅函數，SwiftUI 入面無處不在
  \item 優先 struct，遵守協議（Protocol）；SwiftUI 入面幾乎所有嘢都係 struct
  \item 20\% 語法做 80\% 嘅嘢：先用得，再精通
\end{itemize}

\end{tcbsummary}
